\DocumentMetadata{
  pdfversion=2.0,
  pdfstandard=ua-2,
  lang=en-US,
  tagging=on,
  tagging-setup={math/setup=mathml-SE,tabsorder=structure}
}
\documentclass[11pt,letterpaper]{article}
\usepackage[margin=0.68in,includefoot]{geometry}
\usepackage{newcomputermodern}
\usepackage{amsmath,amscd,mathtools}
\usepackage{tikz,adjustbox}
\providecommand{\square}{\mdlgwhtsquare}
\usepackage[hidelinks]{hyperref}
\hypersetup{
  pdftitle={Analysis First-Year Exam, May 2019, Part 2},
  pdfauthor={University of Florida Department of Mathematics},
  pdfsubject={Graduate mathematics examination},
  pdfdisplaydoctitle=true,
  bookmarksopen=true
}
\setcounter{secnumdepth}{0}
\setlength{\parindent}{0pt}
\setlength{\parskip}{0.28em}
\setlength{\emergencystretch}{3em}
\setlength{\leftmargini}{2.15em}
\setlength{\leftmarginii}{2.65em}
\setlength{\leftmarginiii}{3em}
\pagestyle{empty}
\raggedbottom
\allowdisplaybreaks
\NewTaggingSocketPlug{block/list/label}{examproblem}{%
  \tagstructbegin{tag=Lbl}\tagstructend%
  \tagstructbegin{tag=\LBody}%
  \tagstructbegin{tag=\ProblemHeadingTag,title-o={\ProblemHeadingText}}%
  \tagmcbegin{tag=\ProblemHeadingTag,actualtext={\ProblemHeadingText}}%
  #2%
  \tagmcend%
  \tagstructend%
}
\newenvironment{examproblems}[2]{%
  \def\ProblemHeadingTag{#2}%
  \AssignTaggingSocketPlug{block/list/label}{examproblem}%
  \begin{enumerate}%
  \setcounter{enumi}{#1}%
  \renewcommand{\labelenumi}{\textbf{\theenumi.}}%
  \setlength{\topsep}{0.35em}%
  \setlength{\partopsep}{0pt}%
  \setlength{\itemsep}{0.48em}%
  \setlength{\parsep}{0.18em}%
}{\end{enumerate}}
\newenvironment{examsubparts}{%
  \AssignTaggingSocketPlug{block/list/label}{default}%
  \begin{enumerate}%
  \setlength{\topsep}{0.18em}%
  \setlength{\partopsep}{0pt}%
  \setlength{\itemsep}{0.16em}%
  \setlength{\parsep}{0.1em}%
}{\end{enumerate}}
\newenvironment{examinstructions}{%
  \par\begingroup\itshape
}{%
  \par\endgroup\vspace{0.18em}
}
\makeatletter
\renewcommand\subsection{\@startsection{subsection}{2}{0pt}%
  {-1.25ex plus -.25ex minus -.1ex}%
  {0.55ex plus .1ex}%
  {\normalfont\large\bfseries}}
\renewcommand{\@maketitle}{%
  \newpage\null\vskip -1em%
  {\footnotesize\bfseries
    \href{https://www.ufl.edu/}{University of Florida}, \href{https://math.ufl.edu/}{Department of Mathematics}\par}%
  \vskip -0.5em%
  \tagstructbegin{tag=H1,title={Analysis First-Year Exam, May 2019, Part 2}}%
  \tagpdfparaOff%
  \tagmcbegin{tag=H1}%
    {\LARGE\bfseries\href{https://gma.math.ufl.edu/exams/analysis-first-year/}{Analysis First-Year Exam}, May 2019, Part 2\par}%
  \tagmcend%
  \tagpdfparaOn%
  \tagstructend%
  \vskip 0.25em%
  \tagmcbegin{artifact}%
    \hrule height 1pt%
  \tagmcend%
  \vskip 1em%
}
\makeatother
\title{Analysis First-Year Exam, May 2019, Part 2}
\author{}
\date{}
\begin{document}
\maketitle
\thispagestyle{empty}
\begin{examinstructions}
Do exactly 4 problems. Work must be presented in a neat and logical fashion in order to receive credit. Do not leave any gaps. When a theorem is used in a proof it must be precisely stated.
\end{examinstructions}
\begin{examproblems}{0}{H2}
\def\ProblemHeadingText{Problem 1.}
\item
Let \(f_n:[a,b]\to\mathbb{R}\) be a uniformly bounded sequence of Riemann integrable functions. Prove that if \(f_n\to f\) uniformly, then~\(f\) is Riemann integrable.
\def\ProblemHeadingText{Problem 2.}
\item
Let \(f:[0,1]\to\mathbb{R}\) be continuous. Suppose that
\[
\int_0^1 f(x)x^{2019k}\,dx=0
\]
 for all integers \(k\ge 0\). Must~\(f\) be identically~\(0\)? Prove, or give a counterexample.
\def\ProblemHeadingText{Problem 3.}
\item
Let \(f_n:[0,1]\to\mathbb{R}\) be a sequence of measurable functions and suppose there is a function \(g\in L^1[0,1]\) such that \(|f_n|\le g\) for all~\(n\). Consider the functions
\[
F_n(t)=\int_0^t f_n(x)\,dx.
\]
 Prove
\begin{examsubparts}
\item[\textnormal{i)}]
each \(F_n\) is continuous on \([0,1]\)
\item[\textnormal{ii)}]
there is a subsequence \((F_{n_k})\) converging uniformly on \([0,1]\).
\end{examsubparts}
\def\ProblemHeadingText{Problem 4.}
\item
Let \((X,\mathcal{M})\) be a measurable space and \(f_n:X\to\mathbb{R}\), \(n=1,2,3,\ldots\), a sequence of measurable functions. Prove that each of the following subsets of~\(X\) is measurable:
\begin{examsubparts}
\item[\textnormal{a)}]
\(\{x\mid f_n(x)>0\text{ for infinitely many values of }n\}\)
\item[\textnormal{b)}]
\(\{x\mid \text{the sequence }(f_n(x))\text{ is eventually monotone}\}\)
\item[\textnormal{c)}]
\(\left\{x\,\middle|\,\displaystyle\lim_{n\to\infty}n f_n(x)=0\right\}\)
\end{examsubparts}
\def\ProblemHeadingText{Problem 5.}
\item
Let \(f_1\ge f_2\ge f_3\ge\cdots\) be nonnegative measurable functions on a measure space \((X,\mathcal{M},\mu)\), and put \(f=\lim f_n\). Suppose that \(\int f_k\,d\mu<\infty\) for some~\(k\). Prove that
\[
\int f\,d\mu=\lim\int f_n\,d\mu.
\]
 Give a counterexample to show that the conclusion can fail if the finiteness hypothesis is dropped.
\end{examproblems}
\end{document}
